%=======================================================================
%  main_refs.tex
%
%  Annotated guide to the main references cited in the ContEst paper --
%  especially the ones ContEst is compared against.  Each section describes
%  one reference (or a coherent group): its METHOD, the MATHEMATICS it uses,
%  its BENEFITS, and how it COMPARES against ContEst.
%
%  Compiles standalone (\documentclass{article}); a self-contained
%  bibliography is included at the end.
%=======================================================================
\documentclass[11pt]{article}

\usepackage[margin=1in]{geometry}
\usepackage{amsmath,amssymb,amsthm}
\usepackage{mathtools}
\usepackage{booktabs}
\usepackage{array}
\usepackage{enumitem}
\usepackage{hyperref}
\hypersetup{colorlinks=true,linkcolor=blue,citecolor=blue,urlcolor=blue}

\DeclareMathOperator{\tr}{tr}
\DeclareMathOperator{\diag}{diag}
\DeclareMathOperator*{\argmin}{arg\,min}
\newcommand{\R}{\mathbb{R}}
\newcommand{\E}{\mathbb{E}}
\newcommand{\Th}{\theta}
\newcommand{\ContEst}{\textsc{ContEst}}

% run-in headers for the four aspects of each section
\newcommand{\Method}{\smallskip\noindent\textit{Method.}~}
\newcommand{\Maths}{\smallskip\noindent\textit{Mathematics.}~}
\newcommand{\Benefits}{\smallskip\noindent\textit{Benefits.}~}
\newcommand{\Versus}{\smallskip\noindent\textbf{Comparison with \ContEst.}~}

\title{\bf Main references of the ContEst paper:\\
methods, mathematics, and comparison with ContEst}
\author{ContEst supplementary note}
\date{\today}

\begin{document}
\maketitle

\begin{abstract}
This note collects the principal references cited in the ContEst paper, with an
emphasis on the lines of work ContEst is compared against. Each section explains
the referenced method, the mathematics it uses, its benefits, and---most
importantly---how it relates to and differs from ContEst. A summary table at the
end (Table~\ref{tab:summary}) maps every reference class onto the design axes and
gradient mechanism, mirroring the positioning table of the paper.
\end{abstract}

\tableofcontents

%=======================================================================
\section{What ContEst does, and the axes of comparison}\label{sec:contest}
%=======================================================================

\ContEst\ (Control and Estimation Co-Design) poses actuation \emph{and}
sensing/estimation design in a single two-stage stochastic program. The design
$\Th=(\Th_f,\Th_h)$ collects plant/actuator parameters $\Th_f$ (entering the
dynamics $f$) and sensing/estimator parameters $\Th_h$ (entering the output map
$h$ or the noise covariances), and is chosen to minimize
\begin{equation}\label{eq:contest}
\min_{\Th\in\Theta}\ J_{\mathrm{des}}(\Th)+J_{\mathrm{stoch}}(p_0;\Th),
\qquad
J_{\mathrm{stoch}}=\min_{u_{0:N-1}}\tfrac1N\,\E\Big\{\textstyle\sum_k
\ell(x_k,u_k)+\ell^N(x_N)\Big\},
\end{equation}
the second stage being the stochastic optimal-control (dual-control) value,
expressed through the \emph{information state} (filtering density) so the sensing
operator enters the objective explicitly. The inner value is evaluated by a
\emph{convex} problem---two SDPs (or Riccati equations) in the LQG regime, a pair
of game Riccati equations in the robust $H_\infty$ regime, and a convex
information-state MPC in the eKF-linearized nonlinear regime. The design gradient
$\nabla_\Th J_{\mathrm{stoch}}$ is obtained \emph{exactly} from the dual variables
the inner solver already returns, by the envelope theorem---no differentiation
through the optimizer or its KKT system. Discrete sensor allocation is recovered
by a sparse $\ell_1$ relaxation.

We compare each reference on five axes: \textbf{(i)} is the plant/actuation
designed? \textbf{(ii)} is the sensing/output map designed? \textbf{(iii)} is the
\emph{estimator itself} designed (e.g.\ a filter bandwidth, or the noise level
$V(\Th)$), not merely a choice of sensors? \textbf{(iv)} is closed-loop control in
the design objective? \textbf{(v)} how is the design gradient/derivative obtained?
Nonlinear reach is a sixth, secondary axis.

%=======================================================================
\section{LQG control and sensing co-design (nearest neighbor)}\label{sec:tzoumas21}
%=======================================================================
\noindent\emph{Tzoumas, Carlone, Pappas, Jadbabaie, ``LQG Control and Sensing
Co-Design,'' IEEE~TAC 2021 \cite{tzoumas2021lqg}.}

\Method\ This is the closest single reference to \ContEst. It jointly designs the
sensing (which sensors to activate from a menu, under a sensor budget/cost) and
the LQG controller over a finite horizon. Its central structural result is a
\emph{separation}: the optimal LQG control gains are independent of the sensor
selection, so the co-design collapses to choosing the sensor set that minimizes
the sensing-dependent part of the LQG cost, for which a polynomial-time algorithm
with provable sub-optimality guarantees is given.

\Maths\ Finite-horizon LQG; two Riccati recursions---a control Riccati that does
\emph{not} see the sensors, and a Kalman filtering Riccati whose error covariance
depends on the measurement model. The LQG cost is written as a sum of trace terms
$\sum_t \tr(\Theta_t \Sigma_{t})$ of the filtering error covariances; sensor
selection is a cardinality-constrained set-function problem, shown NP-hard in
general, attacked by a greedy/relaxation algorithm with per-instance and
worst-case bounds derived from the (approximate) supermodular structure.

\Benefits\ A principled, provably near-optimal joint treatment of sensing and LQG
control; identifies the separation that makes the problem tractable; scalable
greedy selection with guarantees.

\Versus\ Same spirit---design sensing \emph{and} control together---but \ContEst\
differs on every axis. (1)~\emph{Discrete vs.\ continuous:} they pick a subset
from a fixed sensor \emph{menu}; \ContEst\ tunes \emph{continuous} sensing/estimator
parameters (precisions, placements on a continuous scale, filter bandwidths,
noise levels $V(\Th)$), recovering the discrete case only when needed via an
$\ell_1$ relaxation. (2)~\emph{Actuation:} their plant/actuation is fixed;
\ContEst\ co-designs $\Th_f$ on the same footing. (3)~\emph{Estimator design:}
they select sensors but do not tune a dynamic estimator; \ContEst\ can design the
estimator itself (e.g.\ the PLL bandwidth study). (4)~\emph{Where separation is
used:} they lean on separation at the \emph{design} level (menu selection given a
fixed control law); \ContEst\ uses separation only \emph{within} the inner solve
and couples both axes through the shared $\Th$, so the co-design does not
separate. (5)~\emph{Gradient:} greedy/relaxation vs.\ an \emph{exact} envelope
gradient. (6)~\emph{Regime:} linear/LQG only vs.\ \ContEst's nonlinear (eKF) and
robust ($H_\infty$) extensions. The paper's PLL study is a head-to-head against
exactly the separation-style sequential pipeline that such menu methods induce.

%=======================================================================
\section{Plant--controller coupling, CCD, and MDO}\label{sec:ccd}
%=======================================================================
\noindent\emph{Fathy, Reyer, Papalambros, Ulsoy, ``On the Coupling Between the
Plant and Controller Optimization Problems,'' ACC 2001 \cite{fathy2001coupling};
Herber \& Allison, ``Nested and Simultaneous Solution Strategies\dots,'' JMD 2019
\cite{herber2019nested}; Allison \& Herber, ``Multidisciplinary Design
Optimization of Dynamic Engineering Systems,'' AIAA~J 2014
\cite{allison2014multidisciplinary}.}

\Method\ This lineage founds \emph{control co-design} (CCD). Fathy et al.\ prove
that plant and controller optimizations are \emph{coupled} and that sequential
(plant-then-control) design is generally suboptimal. Herber \& Allison
systematize the two solution strategies---\emph{nested} (an inner controller solve
inside an outer plant search) and \emph{simultaneous} (one optimization over both).
Allison \& Herber frame dynamic-system CCD as multidisciplinary design
optimization (MDO).

\Maths\ Combined plant/control optimization with first-order (KKT) necessary
conditions; the coupling is captured by cross-sensitivity terms
$\partial^2 J/\partial x_p\,\partial x_c$ between plant $x_p$ and control $x_c$
variables---sequential design is stationary for the joint problem only when these
vanish. Nested strategies form a bilevel program; simultaneous strategies a single
NLP, typically via direct transcription of the dynamics.

\Benefits\ Establishes \emph{why} co-design is necessary; provides a taxonomy and
practical solution architectures for dynamic engineering systems.

\Versus\ \ContEst\ adopts this coupling thesis and extends it to the
\emph{estimation} axis: the sequential plant$\to$sensing$\to$estimator$\to$control
pipeline is a feasible but generally non-stationary point of the joint problem.
Architecturally \ContEst\ is a \emph{nested} (MDF-like) scheme, but its inner
discipline is \emph{convex} and its coupling derivative is supplied \emph{exactly
by duality} (the envelope gradient) rather than by a finite-difference sensitivity
sweep or an adjoint through a nonconvex inner solve---the usual bottleneck of
nested CCD. Crucially, the CCD literature keeps sensing/estimation fixed (full-state
feedback or a fixed observer), so the output map never enters the design; \ContEst\
makes it a first-class design variable.

%=======================================================================
\section{Structural and information-architecture LMI co-design}\label{sec:lmi}
%=======================================================================
\noindent\emph{Grigoriadis, Zhu, Skelton, ``Optimal Redesign of Linear Systems,''
ASME~JDSMC 1996 \cite{grigoriadis1998integrate}; Li, de~Oliveira, Skelton,
``Integrating Information Architecture and Control or Estimation Design,'' SICE~JCMSI
2008 \cite{li2008integrating}.}

\Method\ The structural-control (Skelton) school casts integrated
structure/actuator and control design as \emph{linear matrix inequalities} (LMIs).
Grigoriadis et al.\ alternate a control LMI with a plant/actuator-parameter LMI to
redesign a linear system for better closed-loop covariance performance. Li et al.\
push this to the \emph{information architecture}---sensor/actuator gains and
precisions---treated as decision variables jointly with control \emph{or} estimation.

\Maths\ Covariance (output-variance-constrained) control: express the stationary
$H_2$ cost through a Lyapunov LMI in the closed-loop covariance $\Sigma$ and a
linearizing change of variables ($L=K\Sigma$), giving a convex SDP; alternate
between the control-side and plant/architecture-side LMIs (block coordinate
descent). This is exactly the ancestor of \ContEst's control SDP.

\Benefits\ Each subproblem is convex; structural/architecture parameters enter
tractably; mature LMI machinery and solvers apply.

\Versus\ \ContEst's LQG control SDP \emph{is} the Grigoriadis--Skelton
output-variance-constrained program, and its estimation SDP is the dual. Three
differences. (1)~\emph{Alternation vs.\ exact joint gradient:} they alternate
LMIs (coordinate descent), whereas \ContEst\ reads the \emph{exact} design gradient
from the LMI dual by the envelope theorem and takes a joint quasi-Newton step.
(2)~\emph{``and'' vs.\ ``or'':} Li et al.\ integrate architecture with control
\emph{or} estimation; \ContEst\ designs control \emph{and} estimation together
against one stochastic objective, coupling them through $\Th$. (3)~\emph{Regime:}
their LMIs are linear/$H_2$ (or $H_\infty$); \ContEst\ adds the nonlinear eKF-MPC
surrogate and the fixed-$\gamma$ game-Riccati robust route with the identical
envelope gradient.

%=======================================================================
\section{Sensor selection for estimation}\label{sec:sensorsel}
%=======================================================================
\noindent\emph{Joshi \& Boyd, ``Sensor Selection via Convex Optimization,''
IEEE~TSP 2009 \cite{joshi2009sensor}; Tzoumas, Jadbabaie, Pappas, ``Sensor
Placement for Optimal Kalman Filtering,'' ACC 2016 \cite{tzoumas2016sensor};
Zhang, Ayoub, Sundaram, ``Sensor Selection for Kalman Filtering\dots,'' Automatica
2017 \cite{zhang2017sensor}.}

\Method\ Select a subset of sensors that sharpens an estimation criterion, for a
\emph{fixed} plant and no control in the loop. Joshi \& Boyd maximize a
log-determinant information measure by a Boolean relaxation; Tzoumas et al.\ and
Zhang et al.\ minimize Kalman filtering error by greedy algorithms with
(near-)supermodular guarantees, and characterize when such guarantees hold.

\Maths\ Static case: Fisher information / error-covariance functionals,
$\max_z \log\det\big(\sum_i z_i a_i a_i^\top\big)$ with $z\in\{0,1\}^m$ relaxed to
$[0,1]^m$; convex, then rounded. Dynamic case: Kalman Riccati error covariance as a
set function; submodularity/supermodularity of $\log\det$; greedy with
$(1-1/e)$-type bounds (and counterexamples for other metrics, per Zhang et al.).

\Benefits\ Convex or greedy, scalable to many candidate sensors, with optimality
gaps or approximation guarantees; the canonical tools for pure sensor placement.

\Versus\ These are \emph{estimation-only, open-loop, discrete}. They do not design
actuation, do not tune a dynamic estimator, and do not put closed-loop control in
the objective---the sensor set is optimized for a filtering metric, not for
$J_{\mathrm{stoch}}$. \ContEst\ embeds sensing inside the closed-loop stochastic
control value, co-designs actuation, and returns \emph{exact} design gradients; its
estimation cost $\tr(Q\Sigma^\varepsilon)$ is precisely an A-optimal-type criterion,
so these methods are the open-loop special case of the sensing half of \ContEst.
\ContEst's $\ell_1$ sensor allocation is the closed-loop analog of Joshi--Boyd's
Boolean relaxation.

%=======================================================================
\section{Sparsity-promoting sensor/actuator selection}\label{sec:sparse}
%=======================================================================
\noindent\emph{Dhingra, Jovanovi\'c, Luo, ``An ADMM Algorithm for Optimal Sensor
and Actuator Selection,'' CDC 2014 \cite{dhingra2014admm}; Das \& Bhattacharya,
``Sparse Sensing Architecture for Kalman Filtering\dots,'' 2017
\cite{das2017sparse}; Saraf, Bhattacharya, Skelton, ``$H_2$ Optimal Sensing
Architecture\dots,'' ACC 2017 \cite{saraf2017h2}.}

\Method\ Promote \emph{sparse} sensing/actuation by penalizing an $\ell_1$ (or
reweighted-$\ell_1$) surrogate of the cardinality of the active set, minimizing an
$H_2$/Kalman performance measure; solve by ADMM or SDP, then threshold. Das et al.\
and Saraf et al.\ add guaranteed estimation-error bounds.

\Maths\ $H_2$ norm or Kalman error covariance as the smooth objective;
$\lambda\|\cdot\|_1$ (or reweighted $\sum_i w_i|\cdot|_i$) as the sparsity penalty;
ADMM splitting, or LMI/SDP with sparsity constraints; thresholding to a binary
architecture.

\Benefits\ Convex-surrogate handling of the combinatorial selection; tunable
sparsity--performance trade-off; error guarantees in the guaranteed-bound variants.

\Versus\ This family is the direct ancestor of \ContEst's sensor-allocation
mechanism (Section~7 of the paper): continuous gains $\alpha_j\ge0$ scaling each
sensor, penalized by $\lambda_s\|\alpha\|_1$ in $J_{\mathrm{des}}$, with an epigraph
lift for smoothness. The differences: they optimize for a \emph{fixed} plant/control
and a static or $H_2$ criterion; \ContEst\ places the same $\ell_1$ in a
\emph{closed-loop} (possibly nonlinear) co-design that also moves actuation and
estimator parameters, and keeps the envelope gradient untouched by the penalty.

%=======================================================================
\section{Dual control and stochastic/adaptive MPC}\label{sec:dual}
%=======================================================================
\noindent\emph{Feldbaum, ``Dual Control Theory,'' 1960 \cite{feldbaum1960dual};
Bar-Shalom \& Tse, ``Dual Effect, Certainty Equivalence, and Separation in
Stochastic Control,'' IEEE~TAC 1974 \cite{bar1974dual}; Mesbah, ``Stochastic MPC
with Active Uncertainty Learning: A Survey on Dual Control,'' ARC 2018
\cite{mesbah2018stochastic}.}

\Method\ Dual control recognizes that an input has a \emph{dual effect}: it both
regulates the state and probes to reduce uncertainty about it. Bar-Shalom \& Tse
formalize the dual effect, certainty equivalence (when probing is absent), and
separation. Mesbah surveys how modern stochastic MPC approximates dual control,
distinguishing \emph{active} (probing) from \emph{passive} learning.

\Maths\ Stochastic dynamic programming over the \emph{information state}---the
conditional density $p(x_k\mid\mathcal Z_k)$; the value function depends on this
density, and certainty equivalence holds iff the input has no dual effect. Active
learning augments the cost/constraints to induce informative inputs.

\Benefits\ The correct formulation of control under uncertainty; the conceptual
basis for exploration--exploitation in control.

\Versus\ Dual control is the \emph{objective} \ContEst's second stage adopts:
$J_{\mathrm{stoch}}$ in \eqref{eq:contest} is the dual-control value, and expressing
it through the information state is what lets the sensing operator enter the design.
The distinction is \emph{offline hardware vs.\ online policy}: dual/adaptive MPC
adapts the \emph{inputs} for a \emph{fixed} design $\Th$; \ContEst\ optimizes the
\emph{design} $\Th$ against the (approximated) closed-loop information dynamics. The
two are complementary---the receding-horizon controller keeps running while
\ContEst\ re-provisions parameters. \ContEst\ is honest about the gap: its tractable
eKF surrogate is \emph{caution-aware} (it prices $\tr(Q\Sigma_{t\mid t})$) but,
after linearizing the information dynamics, certainty-equivalent-in-mean---it
captures caution but not active probing, in the Bar-Shalom--Tse sense.

%=======================================================================
\section{Differentiable optimization layers}\label{sec:diffopt}
%=======================================================================
\noindent\emph{Amos \& Kolter, ``OptNet: Differentiable Optimization as a Layer,''
ICML 2017 \cite{amos2017optnet}; Agrawal et al., ``Differentiable Convex
Optimization Layers,'' NeurIPS 2019 \cite{agrawal2019differentiable}.}

\Method\ Embed an optimization problem as a differentiable layer and back-propagate
through its \emph{solution} (the minimizer). OptNet differentiates a QP by
implicitly differentiating its KKT conditions; Agrawal et al.\ generalize to any
convex program by canonicalizing to a cone program and differentiating the
cone-solution (residual) map.

\Maths\ Implicit function theorem applied to the KKT/cone optimality system:
$\partial z^\star/\partial\Th$ is obtained by solving an $m\times m$ linear system
(the differentiated KKT matrix), where $m$ is the number of primal--dual unknowns;
cost $O(m^3)$ per direction. This yields the sensitivity of the \emph{argmin}.

\Benefits\ Fully general end-to-end differentiation of embedded optimizers; ideal
when the downstream loss depends on the \emph{minimizer} $z^\star(\Th)$ itself
(e.g.\ a learned policy or control action).

\Versus\ This is the mechanism \ContEst\ deliberately \emph{avoids}. \ContEst\ needs
the gradient of the inner \emph{value} $V(\Th)$, not of the minimizer; the envelope
theorem supplies it directly from the already-returned dual---\emph{no}
$\partial z^\star/\partial\Th$, \emph{no} KKT-system solve. The complexity gap is
quantitative: implicit differentiation solves an $m\times m$ system with
$m=O(r_x^2)$ for the covariance SDPs, i.e.\ $O(r_x^6)$ per directional derivative,
whereas the envelope contraction is $O(r_x^2 r_\theta)$ for the \emph{whole}
gradient. \ContEst\ uses these tools' theoretical backdrop (sensitivity analysis)
but not their computation.

%=======================================================================
\section{The engine: envelope theorem and perturbation analysis}\label{sec:envelope}
%=======================================================================
\noindent\emph{Milgrom \& Segal, ``Envelope Theorems for Arbitrary Choice Sets,''
Econometrica 2002 \cite{milgrom2002envelope}; Bonnans \& Shapiro, \emph{Perturbation
Analysis of Optimization Problems}, Springer 2000 \cite{bonnans2000perturbation}.}

\Method\ These are \emph{tools}, not competitors. The envelope theorem gives the
derivative of an optimal-value function with respect to parameters as the partial
derivative of the Lagrangian at the optimizer, holding the optimizer fixed.
Bonnans \& Shapiro give the rigorous conditions (Slater, unique primal,
nondegeneracy/strict complementarity) and the directional-derivative formulas for
optimal values of (conic/semidefinite) programs under data perturbations.

\Maths\ For $V(\Th)=\min_z\{\phi(z,\Th):g=0,\,G\in\mathcal K\}$ with optimal triple
$(z^\star,\mu^\star,S^\star)$,
$\nabla_\Th V=\partial_\Th\phi+\langle\mu^\star,\partial_\Th g\rangle
+\langle S^\star,\partial_\Th G\rangle$; the directional derivative equals a
minimization over the dual optimal set, collapsing to a gradient when that set is a
singleton (or, as ContEst shows, when the value is $C^1$).

\Benefits\ Exact value sensitivity computed from the optimizer alone, at negligible
cost beyond the primal--dual solve.

\Versus\ \ContEst's core contribution is recognizing that the control/estimation
\emph{co-design coupling derivative is exactly an envelope gradient}, readable from
the inner SDP, Riccati, or MPC duals. Where the differentiable-optimization layers of
Section~\ref{sec:diffopt} differentiate the minimizer, \ContEst\ invokes these
results to differentiate only the value---turning the coupling derivative from an
expensive sweep into a dual contraction.

%=======================================================================
\section{Adjacent threads used for positioning}\label{sec:adjacent}
%=======================================================================

\noindent\emph{Gevers, ``Identification for Control / Least-Costly Experiment
Design,'' EJC 2005 \cite{gevers2005identification}.}
\Method\ Design the identification experiment of minimum cost subject to a
closed-loop performance requirement. \Versus\ This is the identification-community
sibling of designing the noise level $V(\Th)$ against closed-loop cost; \ContEst\
does the analogous trade \emph{inside} the control loop rather than as a separate
identification stage.

\smallskip
\noindent\emph{Siami \& Jadbabaie, ``A Separation Theorem for Joint Sensor and
Actuator Scheduling,'' Automatica 2020 \cite{siami2021separation}.}
\Method\ Establishes a separation principle for joint sensor/actuator
\emph{scheduling} with guaranteed performance bounds. \Versus\ Like Tzoumas 2021 it
relies on a design-level separation and a discrete (scheduling) decision; \ContEst\
keeps separation only inside the inner solve and treats continuous, coupled design
axes.

\smallskip
\noindent\emph{Zardini, Censi, Frazzoli, ``Co-Design of Autonomous Systems,'' ECC
2021 \cite{zardini2021codesign}.}
\Method\ A compositional (monotone-theory) framework for co-design from hardware
selection to control synthesis. \Versus\ Complementary in scope: it composes
design problems at a systems level, whereas \ContEst\ supplies the continuous,
gradient-based inner co-design of actuation and estimation that such a framework
could call as a component.

\smallskip
\noindent\emph{Ramadan \& Anitescu, ``eKF--Koopman for Tractable Stochastic Optimal
Control,'' IEEE~L-CSS 2024 \cite{ramadan2024extended}; Anderson \& Moore,
\emph{Optimal Filtering}, 1979 \cite{anderson2012optimal}.}
\Method\ The eKF and Koopman-lifted information-state machinery underlying
\ContEst's nonlinear regime. \Versus\ These provide the tractable information state
(and a sharper lift than the plain eKF) that \ContEst's nonlinear surrogate is built
on; the eKF--Koopman lift narrows the moment-closure gap of the surrogate.

%=======================================================================
\section{Summary}\label{sec:summary}
%=======================================================================

Table~\ref{tab:summary} places each reference class against the design axes and the
gradient mechanism. The recurring pattern: prior work designs \emph{one} of
plant/actuation, sensing, or the estimator---and typically differentiates by finite
differences, greedy/relaxation, alternating LMIs, or implicit KKT differentiation.
\ContEst\ closes all three loops in one stochastic objective and obtains the design
gradient exactly from inner duals via the envelope theorem.

\begin{table}[h]
\centering
\caption{Reference classes vs.\ design axes and gradient mechanism.
\checkmark~present, --~absent, $\sim$~partial/menu-only.}
\label{tab:summary}
\small
\renewcommand{\arraystretch}{1.2}
\begin{tabular}{@{}p{4.6cm}ccccp{3.2cm}@{}}
\toprule
reference class & plant/ & sensing & estimator & closed & gradient / \\
                & actuation &        & itself    & loop   & derivative \\
\midrule
CCD / MDO \cite{fathy2001coupling,herber2019nested,allison2014multidisciplinary}
  & \checkmark & -- & -- & \checkmark & finite diff.\ / adjoint \\
LMI structural / info-arch.\ \cite{grigoriadis1998integrate,li2008integrating}
  & \checkmark & $\sim$ & -- & \checkmark & alternating LMIs \\
Sensor selection \cite{joshi2009sensor,tzoumas2016sensor,zhang2017sensor}
  & -- & \checkmark & -- & -- & greedy / relaxation \\
Sparse sensing \cite{dhingra2014admm,das2017sparse,saraf2017h2}
  & $\sim$ & \checkmark & -- & $\sim$ & $\ell_1$ / ADMM \\
LQG sensing co-design \cite{tzoumas2021lqg,siami2021separation}
  & -- & \checkmark\,(menu) & -- & \checkmark & greedy + separation \\
Dual / adaptive MPC \cite{feldbaum1960dual,bar1974dual,mesbah2018stochastic}
  & -- & -- & -- & \checkmark\,(online) & n/a (input policy) \\
Diff.\ opt.\ layers \cite{amos2017optnet,agrawal2019differentiable}
  & -- & -- & -- & \checkmark & implicit KKT diff. \\
\midrule
\textbf{\ContEst}
  & \checkmark & \checkmark\,(cont.+$\ell_1$) & \checkmark & \checkmark
  & exact envelope (dual) \\
\bottomrule
\end{tabular}
\end{table}

%=======================================================================
\begin{thebibliography}{99}
%=======================================================================

\bibitem{tzoumas2021lqg}
V.~Tzoumas, L.~Carlone, G.~J.~Pappas, and A.~Jadbabaie,
\emph{LQG control and sensing co-design},
IEEE Trans.\ Autom.\ Control \textbf{66}(4) (2021), 1468--1483.

\bibitem{fathy2001coupling}
H.~K.~Fathy, J.~A.~Reyer, P.~Y.~Papalambros, and A.~G.~Ulsoy,
\emph{On the coupling between the plant and controller optimization problems},
Proc.\ American Control Conf.\ (ACC), 2001, pp.~1864--1869.

\bibitem{herber2019nested}
D.~R.~Herber and J.~T.~Allison,
\emph{Nested and simultaneous solution strategies for general combined plant and
control design problems},
J.\ Mech.\ Des.\ \textbf{141}(1) (2019), 011402.

\bibitem{allison2014multidisciplinary}
J.~T.~Allison and D.~R.~Herber,
\emph{Multidisciplinary design optimization of dynamic engineering systems},
AIAA Journal \textbf{52}(4) (2014), 691--710.

\bibitem{grigoriadis1998integrate}
K.~M.~Grigoriadis, G.~Zhu, and R.~E.~Skelton,
\emph{Optimal redesign of linear systems},
ASME J.\ Dyn.\ Syst.\ Meas.\ Control \textbf{118}(3) (1996), 598--605.

\bibitem{li2008integrating}
F.~Li, M.~C.~de~Oliveira, and R.~E.~Skelton,
\emph{Integrating information architecture and control or estimation design},
SICE J.\ Control Meas.\ Syst.\ Integr.\ \textbf{1}(2) (2008), 120--128.

\bibitem{joshi2009sensor}
S.~Joshi and S.~Boyd,
\emph{Sensor selection via convex optimization},
IEEE Trans.\ Signal Process.\ \textbf{57}(2) (2009), 451--462.

\bibitem{tzoumas2016sensor}
V.~Tzoumas, A.~Jadbabaie, and G.~J.~Pappas,
\emph{Sensor placement for optimal Kalman filtering: fundamental limits,
submodularity, and algorithms},
Proc.\ ACC, 2016, pp.~191--196.

\bibitem{zhang2017sensor}
H.~Zhang, R.~Ayoub, and S.~Sundaram,
\emph{Sensor selection for Kalman filtering of linear dynamical systems:
complexity, limitations and greedy algorithms},
Automatica \textbf{78} (2017), 202--210.

\bibitem{dhingra2014admm}
N.~K.~Dhingra, M.~R.~Jovanovi\'c, and Z.-Q.~Luo,
\emph{An ADMM algorithm for optimal sensor and actuator selection},
Proc.\ IEEE Conf.\ Decision and Control (CDC), 2014, pp.~4039--4044.

\bibitem{das2017sparse}
N.~Das and R.~Bhattacharya,
\emph{Sparse sensing architecture for Kalman filtering with guaranteed error
bound}, Proc.\ IAA Conf.\ Space Situational Awareness (ICSSA), 2017.

\bibitem{saraf2017h2}
R.~Saraf, R.~Bhattacharya, and R.~E.~Skelton,
\emph{$H_2$ optimal sensing architecture with model uncertainty},
Proc.\ ACC, 2017.

\bibitem{feldbaum1960dual}
A.~A.~Feldbaum,
\emph{Dual control theory. I},
Avtomatika i Telemekhanika \textbf{21}(9) (1960), 1240--1249.

\bibitem{bar1974dual}
Y.~Bar-Shalom and E.~Tse,
\emph{Dual effect, certainty equivalence, and separation in stochastic control},
IEEE Trans.\ Autom.\ Control \textbf{19}(5) (1974), 494--500.

\bibitem{mesbah2018stochastic}
A.~Mesbah,
\emph{Stochastic model predictive control with active uncertainty learning: a
survey on dual control},
Annual Reviews in Control \textbf{45} (2018), 107--117.

\bibitem{amos2017optnet}
B.~Amos and J.~Z.~Kolter,
\emph{OptNet: differentiable optimization as a layer in neural networks},
Proc.\ Int.\ Conf.\ Machine Learning (ICML), 2017.

\bibitem{agrawal2019differentiable}
A.~Agrawal, B.~Amos, S.~Barratt, S.~Boyd, S.~Diamond, and J.~Z.~Kolter,
\emph{Differentiable convex optimization layers},
Advances in Neural Information Processing Systems (NeurIPS), 2019.

\bibitem{milgrom2002envelope}
P.~Milgrom and I.~Segal,
\emph{Envelope theorems for arbitrary choice sets},
Econometrica \textbf{70}(2) (2002), 583--601.

\bibitem{bonnans2000perturbation}
J.~F.~Bonnans and A.~Shapiro,
\emph{Perturbation Analysis of Optimization Problems},
Springer, New York, 2000.

\bibitem{gevers2005identification}
M.~Gevers,
\emph{Identification for control: from the early achievements to the revival of
experiment design},
European J.\ Control \textbf{11}(4--5) (2005), 335--352.

\bibitem{siami2021separation}
M.~Siami and A.~Jadbabaie,
\emph{A separation theorem for joint sensor and actuator scheduling with
guaranteed performance bounds},
Automatica \textbf{119} (2020), 109054.

\bibitem{zardini2021codesign}
G.~Zardini, A.~Censi, and E.~Frazzoli,
\emph{Co-design of autonomous systems: from hardware selection to control
synthesis},
Proc.\ European Control Conf.\ (ECC), 2021, pp.~682--689.

\bibitem{ramadan2024extended}
M.~S.~Ramadan and M.~Anitescu,
\emph{Extended Kalman filter--Koopman operator for tractable stochastic optimal
control},
IEEE Control Systems Letters \textbf{8} (2024).

\bibitem{anderson2012optimal}
B.~D.~O.~Anderson and J.~B.~Moore,
\emph{Optimal Filtering},
Prentice Hall, 1979 (Dover reprint, 2005).

\end{thebibliography}

\end{document}
